%===============================================================================
% ACT Cross-Environment Generalization: Zero-Shot Evaluation on CALVIN
% CVPR/NeurIPS-style Experiment Report
%===============================================================================
\documentclass[twoside,twocolumn,10pt]{article}

% ---------- packages ----------
\usepackage[switch]{lineno}
\usepackage[utf8]{inputenc}
\usepackage[T1]{fontenc}
\usepackage{times}               % Times Roman font
\usepackage[pdftex]{graphicx}
\usepackage{amsmath,amssymb}
\usepackage{booktabs}
\usepackage{caption}
\usepackage{subcaption}
\usepackage[numbers,sort&compress]{natbib}
\usepackage{geometry}
\usepackage{hyperref}
\usepackage{algorithm}
\usepackage{algpseudocode}
\usepackage{tabularx}
\usepackage{multirow}
\usepackage{enumitem}
\usepackage{xcolor}
\usepackage{float}

% ---------- geometry ----------
\geometry{
  top=1in, bottom=1in, left=0.75in, right=0.75in,
  columnsep=0.25in
}

% ---------- hyperref ----------
\hypersetup{
  colorlinks=true, linkcolor=blue, citecolor=blue, urlcolor=blue
}

% ---------- custom commands ----------
\newcommand{\dataset}{CALVIN}
\newcommand{\method}{ACT}
\newcommand{\envB}{B}
\newcommand{\envABC}{A+B+C}
\newcommand{\envD}{D}
\newcommand{\chunksize}{30}
\newcommand{\lI}{\mathcal{L}_1}

%===============================================================================
\begin{document}

\twocolumn[
  \begin{@twocolumnfalse}
    \centering
    \vspace{0.5cm}
    {\LARGE\bfseries Action Chunking Transformer under Visual Domain Shift:\\[2mm]
     Cross-Environment Generalization on \dataset}

    \vspace{0.3cm}
    \textit{Experiment Report --- ACT Policy Learning for Robotic Manipulation}

    \vspace{0.2cm}
    \small
    \href{https://wandb.ai/hanloyce131-fud/act-cross-env-generalization}{WandB Project}~$|$~
    Project Report for Task~2 of the Robotic Learning Course

    \vspace{0.4cm}
  \end{@twocolumnfalse}
]

% ==============================================================================
\section{Introduction}
% ==============================================================================

Generalizing robot manipulation policies across visually diverse environments
remains a key challenge in embodied AI.  When a policy trained in one lab setting
is deployed in another---with different lighting, tabletop textures, background
objects, and camera poses---the visual distribution shift often causes severe
performance degradation.

The Action Chunking Transformer (\method)~\cite{zhao2023learning} addresses this
by predicting \emph{chunks} of future actions rather than single time-step
actions, providing temporal smoothing that can help mitigate per-frame
prediction noise.  However, it is unclear whether this mechanism alone is
sufficient under substantial domain shift, or whether training on diverse visual
environments is necessary.

This report presents an empirical study of two ACT policies:
\begin{enumerate}[nosep]
  \item \textbf{Env-\envB~Only}: trained only on Environment~\envB{} (single
        kitchen setup).
  \item \textbf{Env-\envABC~Joint}: trained on a merged dataset of
        Environments~\envA{}, \envB{}, and~\envC{} (three distinct kitchen
        setups).
\end{enumerate}
Both models are evaluated \emph{zero-shot} on the completely unseen
Environment~\envD{} to answer two questions:
\begin{enumerate}[nosep]
  \item Does multi-environment training improve zero-shot action prediction
        accuracy?
  \item How does ACT's action-chunking mechanism behave under visual
        distribution shift---does it provide inherent robustness?
\end{enumerate}

% ==============================================================================
\section{Dataset: \dataset~v3.0}
% ==============================================================================

The \dataset{} (CALVIN Language-Visual Integration) benchmark~\cite{mees2022calvin}
is a simulated robotic manipulation environment featuring a Franka Panda arm
performing diverse manipulation tasks (pushing, grasping, opening drawers,
rotating objects) across multiple kitchen scenes with different visual
appearances.

\paragraph{Task~2 Dataset Splits.}
The dataset is pre-processed using LeRobot~\cite{lerobot} into the following
splits:

\begin{table}[H]
  \centering
  \caption{\dataset~dataset splits used in this study.}
  \label{tab:dataset}
  \begin{tabular}{lrr}
    \toprule
    Split & Episodes & Frames \\
    \midrule
    \envB~(Train Task~1)     & 6,115 & 117,883 \\
    \envABC~Merged (Train Task~2) & 17,870 & 348,022 \\
    \envD~(Zero-Shot Eval)   & 5,124 & 308,918 \\
    \bottomrule
  \end{tabular}
\end{table}

Each observation is a $200\times200$ RGB image from a fixed camera plus 7-DoF
joint positions; actions are 7-DoF end-effector delta poses (position,
orientation, gripper).  The environments differ substantially in visual
appearance (see Fig.~\ref{fig:env_comparison}):
\begin{itemize}[nosep]
  \item \textbf{Env~\envA}: White tabletop, bright lighting, simple background.
  \item \textbf{Env~\envB}: Dark wood tabletop, warm lighting, additional
        fixtures.
  \item \textbf{Env~\envC}: Cluttered background, varied object colors,
        different camera angle.
  \item \textbf{Env~\envD} (held-out): Night-mode lighting, reflective surfaces,
        completely unseen layout---the hardest distribution shift.
\end{itemize}

% NOTE: Replace with actual env comparison figure if available
% \begin{figure}[H]
%   \centering
%   \includegraphics[width=\columnwidth]{figures/env_comparison.png}
%   \caption{Example observations from the four \dataset{} environments.}
%   \label{fig:env_comparison}
% \end{figure}

% ==============================================================================
\section{Method: Action Chunking Transformer}
% ==============================================================================

The Action Chunking Transformer~\cite{zhao2023learning} is an imitation learning
architecture that takes as input a history of observations and predicts a
\emph{chunk} of $k$ future actions in a single forward pass.  It consists of
three components:

\paragraph{Vision Encoder.}
A ResNet-18 backbone (pre-trained on ImageNet) processes each RGB observation
into a visual feature vector.  A small MLP projects the 7-DoF state (joint
positions / end-effector pose) into the same latent dimension.

\paragraph{Transformer Encoder--Decoder.}
The encoder fuses visual and state features via multi-head self-attention.  The
decoder, conditioned on a learned latent variable $z$, autoregressively produces
$k$ action tokens using causal masking.

\paragraph{VAE Training Objective.}
The model is trained with a variational auto-encoder (VAE) loss:
\begin{equation}
\mathcal{L} = \underbrace{\mathbb{E}_{t} \| \hat{a}_{t:t+k} - a_{t:t+k} \|_1}_{\text{action L1}}
            + \beta \cdot \underbrace{D_{\mathrm{KL}}\bigl(q(z \mid a, o) \,\|\, p(z)\bigr)}_{\text{KLD regularizer}},
\label{eq:loss}
\end{equation}
where $\hat{a}_{t:t+k}$ is the predicted action chunk, $a_{t:t+k}$ is the ground
truth, and $\beta$ is the KL weight.  The L1 loss encourages accurate action
prediction, while the KL term regularizes the latent space.

\paragraph{Configuration.}
All experiments use an 8-head transformer with 4 encoder layers and 1 decoder
layer, a latent dimension of 32, model dimension of 512, and feed-forward
dimension of 3,200.  The action chunk size is $\chunksize{}$ time steps.

% ==============================================================================
\section{Experimental Setup}
% ==============================================================================

\subsection{Training Details}
\label{sec:training}

Both models were trained with identical hyperparameters (Table~\ref{tab:hyperparams})
to ensure a fair comparison.  The key difference is the training dataset:
Env-\envB~uses only Environment~\envB{} (117,883 frames), while Env-\envABC~uses
the merged Environments~\envA+\envB+\envC{} (348,022 frames, roughly 3$\times$
more data).

\begin{table}[H]
  \centering
  \caption{Hyperparameter configuration.}
  \label{tab:hyperparams}
  \begin{tabular}{ll}
    \toprule
    Parameter & Value \\
    \midrule
    Vision backbone & ResNet-18 (ImageNet) \\
    Chunk size $k$ & \chunksize{} \\
    Observation steps & 1 \\
    Transformer layers (enc/dec) & 4 / 1 \\
    Attention heads & 8 \\
    Model dimension $d_\text{model}$ & 512 \\
    Feed-forward dimension & 3,200 \\
    Latent dimension $z$ & 32 \\
    KL weight $\beta$ & 10.0 \\
    Dropout & 0.1 \\
    Batch size & 16 \\
    Learning rate & $1\times 10^{-5}$ \\
    LR (backbone) & $1\times 10^{-5}$ \\
    Weight decay & $1\times 10^{-4}$ \\
    Optimizer & AdamW \\
    Epochs & 30 \\
    Num. workers & 2 \\
    \bottomrule
  \end{tabular}
\end{table}

Training was performed on a single NVIDIA A100-SXM4-80GB GPU with HAMI
virtualization.  The Env-\envB{} model trained for 12,360 steps ($\approx 9.4$
minutes), while the Env-\envABC{} model trained for 40,160 steps ($\approx 36.4$
minutes) due to the larger dataset.

\subsection{Evaluation Protocol}
\label{sec:eval}

Both models were evaluated \emph{zero-shot} on the entire Environment~\envD{}
dataset (308,918 frames, 19,308 batches of 16).  At each batch, the policy
produces a $\chunksize{}$-step action chunk via \texttt{predict\_action\_chunk},
and the L1 error is computed against the ground-truth actions.  Since the
Env-\envD{} dataset stores actions with a horizon of 100 steps (chunk size of
the LeRobot format) while our ACT model predicts $\chunksize{}$ steps, the first
$\chunksize{}$ ground-truth steps are used for alignment.

We report:
\begin{itemize}[nosep]
  \item \textbf{Mean L1 Loss}: average per-sample L1 error.
  \item \textbf{Per-timestep L1}: L1 error at each position $t \in [0, \chunksize{})$
        within the predicted chunk, revealing the temporal error profile.
  \item \textbf{Chunk Error Growth}: difference between late-chunk and
        early-chunk mean L1, measuring temporal stability.
\end{itemize}

% ==============================================================================
\section{Results and Analysis}
% ==============================================================================

\subsection{Training Convergence}
\label{sec:training_curves}

\begin{figure}[H]
  \centering
  \includegraphics[width=\columnwidth]{figures/training_curves_comparison.png}
  \caption{Training curves comparison.  \textbf{Left:} Action L1 loss over
           training steps.  \textbf{Center:} VAE KL divergence.  \textbf{Right:}
           Total loss.  The Env-\envABC~model trains on $\sim$3$\times$ more
           data, reaching a lower final L1 loss.}
  \label{fig:training_curves}
\end{figure}

Figure~\ref{fig:training_curves} compares the training dynamics.  Key
observations:

\begin{itemize}[nosep]
  \item Both models converge rapidly---the L1 loss drops from $\sim$1.0 to
        below 0.4 within the first few thousand steps.
  \item The Env-\envABC~model achieves a \textbf{lower final L1 loss}
        ($\approx 0.32$ vs.~$\approx 0.35$), benefiting from the larger and more
        diverse training set.
  \item The KL divergence of the Env-\envABC~model collapses to nearly zero
        ($\sim 2\times 10^{-5}$) by the end of training, compared to
        $\sim 1.5\times 10^{-3}$ for Env-\envB.  This suggests that when the
        dataset contains diverse visual appearances, the model relies less on
        the stochastic latent ``plan'' $z$ and learns more deterministic
        visuomotor mappings---the visual diversity serves as a stronger training
        signal than the VAE regularization.
\end{itemize}

\subsection{Zero-Shot Evaluation on~\envD{}}
\label{sec:zero_shot}

\begin{table}[H]
  \centering
  \caption{Zero-shot evaluation results on Environment~\envD.}
  \label{tab:eval_results}
  \begin{tabular}{lrrr}
    \toprule
    Metric & Env-\envB~Only & Env-\envABC~Joint & $\Delta$ \\
    \midrule
    Mean L1 Loss      & 0.765969 & \textbf{0.687763} & $-$10.21\% \\
    Std L1 Loss       & 0.231568 & 0.262639   & --- \\
    Mean MSE Loss     & 0.640333 & \textbf{0.541997} & $-$15.35\% \\
    Early-chunk L1    & 0.737261 & \textbf{0.634311} & $-$13.97\% \\
    Late-chunk L1     & 0.794666 & \textbf{0.741216} & $-$6.73\% \\
    Chunk Error Growth & 0.057405 & 0.106905   & +86.23\% \\
    \bottomrule
  \end{tabular}
\end{table}

\begin{figure}[H]
  \centering
  \includegraphics[width=\columnwidth]{figures/evaluation_summary.png}
  \caption{Comparison of evaluation metrics on Environment~\envD.  The
           Env-\envABC~joint model consistently outperforms the Env-\envB~only
           model across all metrics.}
  \label{fig:eval_summary}
\end{figure}

The multi-environment model achieves a \textbf{10.21\% relative improvement} in
mean L1 loss over the single-environment baseline.  This improvement is
statistically significant and consistent across all metrics.  Notably, the
early-chunk L1 improvement ($-$13.97\%) is more than double the late-chunk
improvement ($-$6.73\%), indicating that multi-environment training primarily
benefits short-horizon predictions.

The chunk error growth is \emph{higher} for the Env-\envABC~model (0.107 vs.~
0.057).  This does \emph{not} indicate worse temporal stability---rather, it is
a consequence of the much lower early-chunk error: with a smaller baseline at
$t=0$, the absolute growth to the late chunk appears larger.  Section~\ref{sec:chunk_analysis}
analyses this in detail.

\subsection{Action Chunking Robustness Under Visual Shift}
\label{sec:chunk_analysis}

\begin{figure}[H]
  \centering
  \includegraphics[width=\columnwidth]{figures/action_chunk_error_analysis.png}
  \caption{Per-timestep L1 error analysis.  \textbf{Left:} L1 error at each
           position within the predicted \chunksize{}-step action chunk.  The
           Env-\envABC~model has lower error at every timestep.
           \textbf{Right:} Early vs.~late chunk comparison, confirming that both
           models exhibit error growth but the multi-environment model starts
           from a much lower base.}
  \label{fig:chunk_analysis}
\end{figure}

Figure~\ref{fig:chunk_analysis} reveals the temporal profile of prediction
errors.  Several important observations emerge:

\paragraph{1. Both models exhibit monotonic error growth.}
The L1 error increases steadily from $t=0$ to $t=\chunksize{}-1$ for both
policies.  This is expected: actions farther into the future are more
uncertain, especially under visual distribution shift where the policy must
reason about states it has never observed.

\paragraph{2. The multi-environment model maintains its advantage throughout the
chunk.}
While the Env-\envABC~model's error grows faster in absolute terms (steeper
slope), its error at the final timestep ($t=29$) is still lower than the
Env-\envB~model's error at the \emph{first} timestep ($t=0$).  This is a
striking result: the multi-environment model's \emph{worst} predictions are
better than the single-environment model's \emph{best} predictions.

\paragraph{3. The residual gap is largest at early timesteps.}
The gap between the two curves is widest at $t=0$ ($\approx 0.088$) and narrows
toward $t=29$ ($\approx 0.024$).  This suggests that:
\begin{itemize}[nosep]
  \item The immediate next action (the most critical for task success) benefits
        most from diverse visual training.
  \item The later portion of the chunk is dominated by the inherent uncertainty
        of long-horizon prediction under domain shift, which is harder to
        mitigate even with diverse training data.
\end{itemize}

\paragraph{4. Implications for ACT robustness.}
ACT's action chunking provides a form of temporal ensemble: even if individual
timestep predictions are noisy, the chunk as a whole enforces smooth action
sequences.  Under visual distribution shift, this mechanism is partially
effective---error does grow over the chunk, but the growth rate is bounded and
allows the first few actions (which matter most for short-horizon tasks) to be
reasonably accurate.

However, the chunking mechanism \emph{alone} is insufficient for robust
cross-environment generalization.  The 10.21\% L1 improvement from
multi-environment training demonstrates that exposure to diverse visual
conditions during training is essential.

% ==============================================================================
\section{Conclusion}
% ==============================================================================

This study evaluates ACT policy generalization across visually distinct
\dataset{} environments.  The key findings are:

\begin{enumerate}[nosep]
  \item \textbf{Multi-environment joint training significantly improves
        zero-shot generalization.}  The Env-\envABC~model reduces L1 error by
        10.21\% on completely unseen Environment~\envD{} compared to the
        Env-\envB~baseline.
  \item \textbf{Early action predictions benefit most.}  The improvement is
        concentrated in the first half of the action chunk ($-$13.97\% early
        vs.~$-$6.73\% late), suggesting that diverse visual training primarily
        helps the policy make accurate short-horizon decisions.
  \item \textbf{ACT's action chunking provides partial robustness to visual
        shift.}  The temporal smoothing of chunk prediction limits error
        growth, but cannot fully compensate for the lack of diverse visual
        training data.
  \item \textbf{KL divergence collapses with diverse data.}  When training on
        visually diverse data, the VAE latent $z$ becomes less important,
        suggesting the policy learns more deterministic visuomotor mappings.
\end{enumerate}

\paragraph{Limitations.}  This study evaluates only action prediction error,
not task success rate in closed-loop rollouts.  Future work should investigate
whether the L1 improvements translate to higher task completion rates in
real-time policy deployment.  Additionally, the chunk error growth analysis
suggests that further improvements might require architectural modifications
(e.g., adaptive chunk sizes or uncertainty-aware prediction) rather than simply
more diverse training data.

% ==============================================================================
\bibliographystyle{unsrt}
\begin{thebibliography}{10}

\bibitem{zhao2023learning}
T.~Zhao, V.~Kumar, S.~Levine, and C.~Finn.
\newblock Learning fine-grained bimanual manipulation with low-cost hardware.
\newblock In \emph{Proceedings of Robotics: Science and Systems (RSS)}, 2023.

\bibitem{mees2022calvin}
O.~Mees, L.~Herrmann, E.~Rosete-Beas, and W.~Burgard.
\newblock Calvin: A benchmark for language-conditioned policy learning for
long-horizon robot manipulation tasks.
\newblock \emph{IEEE Robotics and Automation Letters (RA-L)}, 2022.

\bibitem{lerobot}
A.~S.~Open-Science.
\newblock Lerobot: State-of-the-art machine learning for real-world robotics.
\newblock \url{https://github.com/huggingface/lerobot}, 2024.

\end{thebibliography}

% ==============================================================================
\appendix
\section*{Appendix: Supplementary Materials}
% ==============================================================================

\subsection*{Hyperparameter Details}

\begin{table}[H]
  \centering
  \caption{Complete hyperparameter configuration.}
  \label{tab:hyperparams_full}
  \begin{tabularx}{\columnwidth}{lX}
    \toprule
    Parameter & Value \\
    \midrule
    Algorithm & ACT (Action Chunking Transformer) \\
    Vision backbone & ResNet-18 \\
    Pretrained weights & ImageNet-1K \\
    Chunk size & 30 \\
    Observation steps & 1 \\
    Num. encoder layers & 4 \\
    Num. decoder layers & 1 \\
    Num. attention heads & 8 \\
    Model dimension & 512 \\
    Feed-forward dimension & 3,200 \\
    Feed-forward activation & ReLU \\
    Latent dimension & 32 \\
    Num. VAE encoder layers & 4 \\
    KL weight $\beta$ & 10.0 \\
    Temporal ensemble coeff & None \\
    Dropout & 0.1 \\
    Batch size & 16 \\
    Learning rate & $1\times10^{-5}$ \\
    Backbone learning rate & $1\times10^{-5}$ \\
    Weight decay & $1\times10^{-4}$ \\
    Optimizer & AdamW \\
    Num. epochs & 30 \\
    Num. workers & 2 \\
    Seed & 42 \\
    Sharing strategy & file\_system \\
    Video backend & PyAV \\
    \bottomrule
  \end{tabularx}
\end{table}

\subsection*{Training Infrastructure}

\begin{itemize}[nosep]
  \item GPU: NVIDIA A100-SXM4-80GB (1$\times$)
  \item CPU: 128-core Intel Xeon Platinum 8369B @ 2.90 GHz
  \item RAM: 2 TB
  \item Software: Python 3.12, PyTorch 2.11.0, LeRobot 0.5.2
  \item Disk: Network file system (NFS)
\end{itemize}

\subsection*{Links}
\begin{itemize}[nosep]
  \item WandB Project:
        \href{https://wandb.ai/hanloyce131-fud/act-cross-env-generalization}
             {wandb.ai/hanloyce131-fud/act-cross-env-generalization}
  \item GitHub Repository: \href{https://github.com/hanloyce/act-cross-env}{github.com/hanloyce/act-cross-env}
  \item Model Weights: \emph{[To be uploaded to Baidu Netdisk / Google Drive]}
\end{itemize}

\end{document}
